\documentclass[10pt,a4paper]{article}
\usepackage{fontspec,amsmath,amssymb,graphicx}
\usepackage[margin=18mm]{geometry}
\setmainfont{Times New Roman}
\setlength{\parindent}{0pt}
\setlength{\parskip}{5pt}
\allowdisplaybreaks
\sloppy
\emergencystretch=2em
\begin{document}
\section*{BT\_2\_3\_13 — Lời giải bổ sung}
\subsection*{Kiến thức và trình tự giải}
\textbf{Mục tiêu.} Giải phương trình hyperbolic sau khi đã tìm hai tọa độ đặc trưng độc lập.

\textbf{Bước 1 — Phân loại, tìm đặc trưng và thay đạo hàm.} Làm đủ các bước của phép đưa về chính tắc. Viết cả phép đổi ngược để biểu diễn vế phải $\widetilde f(\xi,\eta)=f(x(\xi,\eta),y(\xi,\eta))$.

\textbf{Bước 2 — Tích phân phương trình chính tắc.} Với $v_{\xi\eta}=g$, giữ biến kia cố định khi tích phân:


\[
v=F(\xi)+G(\eta)+\int_{\xi_0}^{\xi}\int_{\eta_0}^{\eta}g(a,b)\,db\,da.
\]


Nếu tìm được một nghiệm riêng $P(x,y)$ dễ hơn, viết $u=P+F(\xi)+G(\eta)$. Với $v_{\xi\eta}+\alpha v_\eta=0$, đặt $w=v_\eta$, giải $w_\xi+\alpha w=0$ rồi tích phân theo $\eta$; thu được $v=F(\xi)+e^{-\alpha\xi}G(\eta)$.

\textbf{Bước 3 — Đổi cả dữ liệu.} Thế tọa độ đường dữ liệu vào $\xi,\eta$, dùng quy tắc dây chuyền cho đạo hàm được cho. Trên đường không đặc trưng, hai dữ liệu thường xác định hai hàm $F,G$. Trên đường đặc trưng, có biến không thay đổi, nên dữ liệu chỉ xác định giá trị hoặc đạo hàm của một hàm tại một điểm; phải tìm điều kiện tương thích và nêu phần tùy ý còn lại.

\textbf{Bước 4 — Lập đẳng thức hình bình hành.} Với $v=F(\xi)+G(\eta)$, bốn đỉnh theo hai họ đặc trưng thỏa $v_{11}+v_{22}=v_{12}+v_{21}$. Nếu có số hạng mũ, trước hết nhân trọng số để được tổng hai hàm một biến. Với vế phải $g$, hiệu bốn đỉnh bằng tích phân kép của $g$ trên hình chữ nhật trong tọa độ mới.

\textbf{Bước 5 — Xác định miền được dữ liệu quyết định.} Hai đặc trưng đi qua điểm đang xét phải gặp đoạn dữ liệu phù hợp. Với phương trình sóng $v_{tt}=v_{zz}$ trên $z>0$, dữ liệu lúc $t=0$ xác định duy nhất miền $z>t$; ở $0<z<t$, cần thêm dữ liệu biên $z=0$. Với đặc trưng cong, dùng miền trong các tọa độ mới rồi đổi lại.

\textbf{Bước 6 — Kiểm tra.} Thế nghiệm vào cả PDE và từng điều kiện. Nếu đề yêu cầu hai nghiệm, chỉ rõ hai lựa chọn khác nhau cho hàm tùy ý và chứng minh phần chênh lệch thỏa phương trình thuần nhất, đồng thời triệt tiêu trên đường dữ liệu.
\subsection*{Lời giải áp dụng cho bài}
\textbf{(a) Phân loại và đổi biến.} $A=(1+x^2)^2$, $B=0$, $C=-1$, $\delta=(1+x^2)^2>0$. Đặt $z=\arctan x\in(-\pi/2,\pi/2)$; $z_x=1/(1+x^2)$, $z_{xx}=-2x/(1+x^2)^2$. Thế vào
\[u_x=\frac{v_z}{1+x^2},\quad u_{xx}=\frac{v_{zz}-2xv_z}{(1+x^2)^2}\]
cho $v_{zz}-v_{yy}=0$: hai hạng $v_z$ triệt tiêu. Chọn $\xi=z+y$, $\eta=z-y$ cho $4v_{\xi\eta}=0$.

\textbf{(b) Nghiệm tổng quát.} $u=F(\arctan x+y)+G(\arctan x-y)$; hai họ đặc trưng $\arctan x\pm y=\mathrm{const}$.

\textbf{(c) Dữ liệu trên đặc trưng $y=z$.} Đặt $\Phi(z)=\varphi(\tan z)$, $\Psi(z)=\psi(\tan z)$. Tại $\eta=0$, $\xi=2z$,
\[F(2z)+G(0)=\Phi(z),\qquad F'(2z)-G'(0)=\Psi(z).\]
Lấy đạo hàm dòng đầu theo $z$ cho $2F'(2z)=\Phi'(z)$. Vậy điều kiện cần và đủ là
\[\boxed{\frac12\Phi'(z)-\Psi(z)=\lambda\quad\text{hằng số},\quad |z|<\pi/2.}\]
Trong biến gốc: $(1+x^2)\varphi'(x)/2-\psi(x)=\lambda$. Khi điều kiện đúng, $F(s)=\Phi(s/2)-G(0)$ ở $|s|<\pi$, còn $G$ tùy ý với $G'(0)=\lambda$. Lấy $G_1(s)=\lambda s$, $G_2(s)=\lambda s+s^2$:
\[u_1=\varphi\!\left(\tan\frac{\arctan x+y}{2}\right)+\lambda(\arctan x-y),\qquad
u_2=u_1+(\arctan x-y)^2.\]
Hai công thức xác định trên miền $|\arctan x+y|<\pi$. Khi $y=\arctan x$, $u_j=\varphi(x)$ và $(u_j)_y=(1+x^2)\varphi'(x)/2-\lambda=\psi(x)$. Trong tọa độ đặc trưng, mỗi nghiệm là tổng hai hàm một biến, nên thỏa PDE. Nếu yêu cầu nghiệm trên toàn $\mathbb R^2$, phải chọn kéo dài $C^2$ của $F$ ra ngoài $(-\pi,\pi)$; dữ liệu chỉ ở $x$ hữu hạn không bảo đảm kéo dài trơn tại hai đầu nếu $\varphi$ tăng không kiểm soát. Vì vậy công thức trên khẳng định tồn tại trên miền tự nhiên chứa toàn đường dữ liệu.\end{document}